\documentclass[]{article}
\usepackage[utf8]{inputenc}
\usepackage[T1]{fontenc}
\usepackage[french]{babel}
\usepackage{graphicx}
\usepackage{geometry}
\geometry{margin=2.5cm}
%opening
\title{\textbf{RESUME DU COURS D'INVESTIGATION NUMERIQUE}}
\date{}

\begin{document}
\thispagestyle{empty}
\begin{center}
	\includegraphics[width=0.5\textwidth]{logo_po}\\
	\vspace{\baselineskip}
	\rule{0.95\textwidth}{2pt}
	\vspace{\baselineskip}\large\textbf{THEORIES ET PRATIQUES DE L’INVESTIGATION  NUMERIQUE.}\normalsize\\
	\rule{0.95\textwidth}{2pt}\vspace{\baselineskip}\vfill
	\vspace{\baselineskip}
	\textbf{Enseignant: M.MINKA THIERRY.}\\
	Par l'étudiante: CHEUTCHEULONGA HILARY PAULA\\
	Filière CIN 4\vfill
	\rule{\textwidth}{0.5pt}
	\vspace{\baselineskip}\emph{\textbf{Année Academique: 2025-2026.}}\vspace{\baselineskip}\\
	\rule{\textwidth}{0.5pt}
\end{center}
\newpage
\begin{center}
	\large\textbf{DEVOIR : CHAPITRE 1}\\
\end{center}
\section*{Partie 1 : Fondements Philosophiques et Épistémologiques}

\subsection*{1. Analyse critique du paradoxe de la transparence}

Byung-Chul Han, dans La Société de la Transparence (2015), met en lumière un paradoxe fondateur de notre époque numérique : la quête obsessionnelle de transparence, censée renforcer la confiance et la démocratie, conduit paradoxalement à la disparition de l’intimité et à l’érosion de la liberté. L’idéologie de la transparence part de l’idée que « voir tout » permet d’éviter la corruption, les abus ou les manipulations. Mais, appliquée sans limites, elle se transforme en mécanisme de surveillance généralisée, où l’individu devient totalement exposé et réduit à un objet de données.
Dans une perspective philosophique, la transparence absolue n’est pas synonyme de vérité. Elle engendre plutôt une logique de contrôle : ce qui est rendu visible devient mesurable, comparable et exploitable. Le paradoxe réside donc dans le fait que la transparence, censée protéger la démocratie et la confiance, en vient à instaurer une société disciplinaire plus proche du panoptique de Bentham que d’un espace libre d’échanges.
À l’ère numérique, ce paradoxe prend une forme accrue. Les plateformes sociales incitent les individus à s’exposer en permanence, réduisant la frontière entre vie privée et vie publique. Or, plus un sujet se dévoile, plus il perd sa capacité de retrait, condition fondamentale de l’autonomie. Comme le note Han, l’excès de transparence supprime le « mystère » et l’altérité, dimensions essentielles de l’être humain.
Dans le domaine de l’investigation numérique, ce paradoxe est particulièrement manifeste. Un gouvernement, par exemple, peut justifier une surveillance de masse au nom de la transparence et de la sécurité nationale. Mais ce processus conduit à une vulnérabilité structurelle : les citoyens perdent la maîtrise de leurs données personnelles, et les traces numériques deviennent une arme de contrôle politique. La quête de transparence se retourne ainsi contre sa finalité démocratique initiale.
-Cas concret : transparence gouvernementale vs vie privée
Un exemple classique est celui des bases de données biométriques nationales. D’un côté, elles permettent une transparence accrue des procédures administratives (authentification rapide, traçabilité des identités). De l’autre, elles mettent en péril la vie privée : centralisation massive de données sensibles, risques de piratage ou d’usage politique abusif. La volonté de transparence de l’État aboutit alors à une surveillance permanente des citoyens, fragilisant les droits fondamentaux.
-Résolution pratique inspirée de l’éthique kantienne
L’éthique kantienne propose un critère simple mais exigeant : agir de telle sorte que la maxime de son action puisse valoir comme loi universelle. Appliqué à l’investigation numérique, cela signifie que les pratiques de transparence ne doivent jamais instrumentaliser l’individu comme simple moyen de contrôle, mais le respecter comme fin en soi.Une résolution kantienne au paradoxe serait d’instaurer une transparence proportionnée et réciproque :
*Les institutions doivent être transparentes envers les citoyens (sur leurs décisions, leurs algorithmes, leurs critères d’usage des données).
*Les citoyens, eux, doivent conserver un droit au retrait, c’est-à-dire la possibilité de ne pas être totalement exposés dans leur intimité numérique.
*La solution réside dans un équilibre : la transparence doit viser la responsabilité institutionnelle, et non la mise à nu totale de l’individu. Kant inviterait ici à formuler un impératif catégorique du numérique : « N’utilise jamais la donnée personnelle uniquement comme moyen de surveillance, mais toujours en respectant la dignité et l’autonomie de la personne. »
\subsection*{2. Transformation ontologique du numérique}
-Comparaison Heidegger / ère numérique
Chez Heidegger, l’être se comprend comme être-au-monde, inscrit dans une temporalité et une matérialité. La technique moderne, selon lui, risque de réduire l’être à un simple « fonds disponible » (Bestand). Dans l’ère numérique, cette réduction s’accentue : l’identité humaine est prolongée et parfois enfermée dans son double numérique. L’individu n’est plus seulement présent par son corps, mais par ses traces, ses profils et ses données.
- Profil social comme « être-par-la-trace »
Un profil Facebook ou LinkedIn est une manifestation d’« être-par-la-trace » : l’existence numérique se construit à travers des publications, des likes, des historiques. Même en l’absence physique de la personne, son « être » reste accessible et interprétable à travers ces traces. L’ontologie se déplace : exister, c’est laisser des données.
-Impact sur la preuve légale
Cette transformation ontologique bouleverse la notion de preuve. La trace numérique, volatile et manipulable, devient centrale dans les enquêtes judiciaires. Mais elle n’a pas la stabilité d’un objet matériel : elle exige des protocoles de conservation (chaîne de confiance, hash, métadonnées). En outre, l’« être-par-la-trace » pose un problème d’authenticité : une identité peut être usurpée, un profil falsifié. Ainsi, la preuve légale doit évoluer vers une épistémologie de la probabilité et de la chaîne de validation, plutôt que de la certitude matérielle.
\section*{Partie 3 : Révolution Quantique et ses Implications}
\subsection*{6. Expérience de pensée : Chat de Schrödinger numérique}
-Conception (version numérique)
Imaginons un fichier numérique F stocké sur un dispositif quantique (ou dans un système hybride classique/quantique) dont l’état informationnel n’est pas encore mesuré/validé. Avant toute observation forensique, on modélise F comme étant en superposition entre deux branches conceptuelles :
présent : le fichier est lisible et intact ;
effacé : le fichier a été effacé/altéré.
L’état global peut s’écrire (idéellement) :
|Ψ⟩ = α|présent⟩ + β|effacé⟩, avec |α|²+|β|² = 1.
-Existe-t-il vraiment une superposition « présent/effacé » avant analyse ?
Pragmatiquement : non pas comme un chat vivant/mort réel, mais oui en modèle épistémique — tant qu’on n’a pas fait d’opération qui décide définitivement de l’état (p.ex. lecture destructrice, restauration complète ou preuve immuable), l’investigateur travaille avec une incertitude quantifiée sur l’état. Sur supports classiques, l’« état superposé » est une métaphore utile pour désigner indétermination due à volatilité, réplication concurrente, ou altérations possibles. Sur supports réellement quantiques (qubit-based storage), la superposition est littérale et la mesure la projettera.
-Impact sur la notion de preuve « certaine » en justice
La certitude absolue (100 %) devient impossible si :
la lecture peut modifier l’objet (effet de la mesure) ;
la copie parfaite est impossible (non-clonage, voir point 8).
On passe d’une preuve certaine à une preuve probabilisée : la chaîne de confiance doit fournir un niveau de confiance (1−ε) avec ε quantifié. Les tribunaux devront accepter des attestations probabilistes assorties de procédures de validation mathématiquement formalisées (marges d’erreur, intervalle de confiance, preuves cryptographiques).
-Protocole d’observation minimisant l’effet (pratique forensique inspirée des QND \& ZK)
*Isolement physique / électronique du support (standard chain-of-custody).
*Création d’un ancilla (registre auxiliaire) qui s’entangle faiblement avec l’état du fichier afin d’extraire une propriété sans projeter complètement (analogue à une mesure faible / quantum non-demolition — QND).
*Mesure sur l’ancilla uniquement pour obtenir des métadonnées certifiables (par ex. existence, taille, empreinte probabiliste) sans détruire l’état primaire.
*Calcul immédiat d’une engagement cryptographique (commitment) post-mesure : un engagement chiffré / hash post-quantique horodaté (ledger distribué ou TSP) signé par l’investigateur (post-quantum signature).
*Production d’une preuve zero-knowledge (ZK) qui atteste certaines propriétés (p.ex. « le fichier contient telle structure ») sans révéler le contenu, couplée à une non-répudiation (ZK-NR).
*Si nécessaire, effectuer une copie classique après transformation contrôlée (mesure projetante) documentée par timestamp et attestations.
*Archive immuable des preuves classiques dérivées (hashes, preuves ZK, logs) et conservation du support originel dans environnement isolé.
*Ce protocole cherche à minimiser l’effet « mesurant = modifiant » en extrayant d’abord des preuves non-destructives et en horodatant/descriptant rigoureusement toute action invasive ultérieure.
\subsection*{7. Calculs sur la sphère de Bloch (état donné)}
$État fourni :|ψ⟩ = cos(π/6)|0⟩ + e^{iπ/4} sin(π/6)|1⟩.
Remarque : cela correspond au formalisme standard où l’angle sur la sphère est θ = π/3 (puisque cos(θ/2)=cos(π/6)\\, 
sin(θ/2)=sin(π/6)) et ϕ = π/4.\\
-Probabilités de mesure P(0) et P(1)\\
cos(π/6) = √3 / 2.\\
√3 ≈ 1.7320508075688772 → √3 / 2 ≈ 0.8660254037844386.\\
sin(π/6) = 1/2 = 0.5.\\
Probabilité P(0) = |cos(π/6)|² = (√3 / 2)² = 3 / 4 = 0.75.\\
Détail numérique : 0.8660254037844386² = 0.7499999999999999 ≈ 0.75.\\
Probabilité P(1) = |sin(π/6)|² = (1/2)² = 1/4 = 0.25.\\
Vérification somme : 0.75 + 0.25 = 1.00.\\
-Représentation sur la sphère de Bloch
Formule (pour θ,ϕ) : vecteur de Bloch = (x,y,z) = (sin θ cos ϕ, sin θ sin ϕ, cos θ).\\
Ici θ = π/3, ϕ = π/4.\\
Calculs :\\
sin(π/3) = √3 / 2 ≈ 0.8660254037844386.\\
cos(π/4) = √2 / 2 ≈ 0.7071067811865476.\\
sin(π/4) = √2 / 2 ≈ 0.7071067811865476.\\
cos(π/3) = 1/2 = 0.5.\\
x = sin θ · cos ϕ ≈ 0.8660254037844386 × 0.7071067811865476.\\
Multiplication : 0.8660254037844386 × 0.7071067811865476 ≈ 0.6123724356957945.\\
y = sin θ · sin ϕ ≈ same product ≈ 0.6123724356957945.\\
z = cos θ = 0.5.\\
=> Bloch vector ≈ (0.6123724, 0.6123724, 0.5).\\
(Graphique) Sur la sphère, ce point est situé dans l’octant où x>0,y>0,z>0, à 60° (π/3) du pôle nord, avec azimut ϕ = 45°. Visuellement : incliné vers la « diagonale » x=y et légèrement au-dessus de l’équateur (z=0.5).
-Impact sur un système de preuve quantique
P(0)=0.75 signifie qu’une mesure dans la base computational fournit l’état |0⟩ avec probabilité 75 %. Une seule mesure ne donne pas certitude ; il faut répéter ou combiner avec preuves complémentaires.
Sur la sphère, l’angle θ/ϕ encode la relation entre la « propriété » qu’on veut prouver et la base de mesure. Choisir une mauvaise base revient à perdre information (basses probabilités).
Dans un système de preuve quantique, les preuves sont statistiques : il faut protocoles (tests de répétition, tomography partielle, échantillonnage) et assurances cryptographiques pour atteindre la confiance requise.
$\subsection*{8. Théorème de non-clonage — conséquences et alternative ZK-NR}
-Pourquoi le théorème empêche la copie parfaite (preuve simple)
Supposons qu’il existe une opération linéaire unitaire U telle que, pour tout état arbitraire |ψ⟩ et un état de référence |0⟩ :
U(|ψ⟩ ⊗ |0⟩) = |ψ⟩ ⊗ |ψ⟩.
Soit aussi |φ⟩ un autre état non orthogonal. Alors :
U(|φ⟩ ⊗ |0⟩) = |φ⟩ ⊗ |φ⟩.
Conservant le produit scalaire (unitarité), on doit avoir :
⟨ψ|φ⟩ = ⟨ψ|φ⟩².
Soit s = ⟨ψ|φ⟩ (complexe). L’égalité s = s² implique s ∈ {0,1}. Donc soit états orthogonaux (s=0) soit identiques (s=1). Contradiction pour états non orthogonaux généraux → pas d’unique U qui clone tous les états. QED.
-Implications pour conservation des preuves quantiques
Impossible de faire des copies parfaites d’un état quantique inconnu pour l’archivage : on ne peut pas créer une « copie de sauvegarde » parfaite sans altérer l’original.
Mesurer pour extraire une copie classique transforme l’état et peut détruire l’information quantique.
Chaîne de conservation doit se reposer sur des attestations cryptographiques et des preuves dérivées (hashs, preuves ZK) plutôt que sur duplications d’état quantique.
-Alternative pratique utilisant ZK-NR (Zero-Knowledge + Non-Repudiation)
Combiner deux idées : (1) extraire preuves classiques vérifiables sans révéler le contenu via ZK ; (2) fournir non-répudiation post-quantique par signatures/horodatages résistants aux attaques quantiques.Protocole candidate (concept) :
*Préparation : si l’objet est quantique, effectuer une opération contrôlée d’entanglement avec ancilla et extraire propriétés vérifiables (p.ex. stabilizers) via mesures non-projectives quand possible.
*Commitment post-mesure : calculer un engagement cryptographique (commitment) des sorties mesurées ; signer avec une signature post-quantique (ex : Dilithium/Falcon si requis) et horodater via ledger distribué.
*Preuve Zero-Knowledge (ZK-NR) : générer une preuve zero-knowledge attestant l’existence et certaines propriétés (par ex. « contient l’empreinte H » ou « propriété P est vraie ») sans révéler l’état. Intégrer une composante non-répudiation (preuve d’origine, identité du déposant/investigateur).
*Archivage : stocker l’engagement, la preuve ZK, la signature, et les logs (audit trail). Si la loi exige exposition complète, on peut ensuite procéder à une mesure projetante contrôlée documentée.
\section*{Partie 5 : Intégration et Synthèse Avancée}
\subsection*{11. Étude de Cas Complexe : Affaire "QuantumLeaks"}
-Cet exercice demande l'élaboration d'un rapport avec des recommandations techniques et éthiques pour un scénario de fuite de documents classifiés chiffrés post-quantique, avec pour défi la conciliation du Trilemme CRO (Confidentialité, Fiabilité (Reliability), et Opposabilité juridique)  pour une preuve devant être préservée pour plus de 30 ans.
*Recommandations Techniques
L'objectif principal est de garantir la non-répudiation post-quantique et la fiabilité de la preuve sur le long terme.Conservation de la Preuve Quantique (Chain of Custody) :
Adoption de Protocoles ZK-NR : Utiliser les Protocoles Zero-Knowledge Non-Repudiation (ZK-NR)  pour toutes les étapes de la chaîne de conservation (Chain of Custody). Ces protocoles sont conçus pour prouver l'authenticité et l'intégrité de la preuve sans en divulguer le contenu, conciliant ainsi authenticité (A(P)) et confidentialité (C(P)).
Métrologie Quantique : Remplacer la "certitude à 100%" par une preuve probabiliste, une "confiance à 1−ϵ avec ϵ<2 −λ" , en utilisant de nouveaux étalons de mesure pour l'incertitude et la fiabilité.
Résistance Post-Quantique (Preuve à 30+ ans) :Migration Cryptographique : S'assurer que tous les mécanismes d'horodatage et de signature numérique appliqués aux preuves sont basés sur des algorithmes de cryptographie post-quantique, tels que Dilithium/Falcon, pour garantir la non-répudiation au-delà de 2030.
Système Q2CSI : Utiliser une Infrastructure de Sécurité Contextuelle et Composable Quantique (Q2CSI)  pour gérer la signature cryptographique de manière légalement explicable et résiliente aux menaces quantiques.
Garantie d'Authenticité (Non-Clonage) :Compte tenu du Théorème de Non-Clonage quantique qui empêche la duplication parfaite d'un état quantique inconnu, il est crucial d'enregistrer et d'auditer chaque mesure/observation de la preuve, car l'acte d'observation lui-même modifie irréversiblement l'état observé (Observateur Participatif).
*Recommandations Éthiques
Le défi éthique consiste à naviguer dans le Trilemme Éthique Fondamental CRO (max{T,E,I} sujet à T+E+I≤2+ϵ)  dans un contexte de sécurité nationale.
Conciliation CRO :Confidentialité (C) / Transparence (T) : Le secret inhérent aux documents classifiés (Confidentialité) doit être mis en balance avec la Transparence démocratique (T). L'utilisation des protocoles ZK-NR permet d'atteindre la Fiabilité et l'Opposabilité sans compromettre complètement la Confidentialité en révélant uniquement le fait de l'attestation, et non le contenu.
Intégrité Procédurale (I) : L'investigateur doit faire preuve d'Intégrité procédurale (I) , garantissant que la quête d'Efficacité investigative (E) pour la sécurité nationale ne mène pas à une altération (même involontaire) de la preuve, respectant ainsi l'impératif éthique de l'Archéologue du digital.
Impératifs Post-Quantiques :Impératif Intergénérationnel : Considérer l'impact de l'investigation sur les générations futures dans un "monde de mémoire parfaite", en s'assurant que les preuves conservées ne pourront pas être compromises par des technologies futures, mais aussi que leur divulgation future respecte un équilibre éthique.
Impératif de l'Oubli : Si des données non pertinentes ou personnelles sont capturées, elles doivent être traitées avec respect pour le droit fondamental à l'oubli numérique.

\end{document}
